\chapter{AutoReg and ARDL}
\label{cap:ardl}
% ── index ───────────────────────────────────────────────────────
\index{ARDL|textbf}
\index{ardl@\texttt{ardl()}|textbf}
\index{AutoReg|textbf}
\index{autoreg@\texttt{autoreg()}|textbf}
\index{distributed lags|textbf}
\index{long-run multiplier|textbf}
\index{Pesaran, cointegration test of}
\index{bounds test|see{Pesaran, cointegration test of}}

% ─────────────────────────────────────────────────────────────────────────────
\section{Models with lags in y and x}

The ARIMA model (ch.~\ref{cap:arima}) deals only with the dynamics of $y$.
In economics, however, the current level of $y$ depends not only on its own past,
but also on current and past values of regressors $x$. The
\emph{Autoregressive Distributed Lag} (ARDL) model captures exactly this.

\subsection*{AutoReg(p): pure AR via OLS}

The AR($p$) model writes $y_t$ as a linear function of its own lags:
\[
  y_t = c + \sum_{i=1}^p \phi_i\, y_{t-i} + \varepsilon_t
\]
Unlike ARIMA (estimated by MLE), \texttt{autoreg()} estimates via
OLS. This is computationally simpler and asymptotically equivalent
for stationary processes. The \texttt{trend} option allows including
a constant and/or a deterministic trend:

\begin{center}
\begin{tabular}{ll}
\toprule
\textbf{trend} & \textbf{Deterministic component} \\
\midrule
\texttt{"c"}  & Constant (default) \\
\texttt{"ct"} & Constant + linear trend \\
\texttt{"t"}  & Trend only \\
\texttt{"n"}  & None \\
\bottomrule
\end{tabular}
\end{center}

\subsection*{ARDL(p, q): distributed lags}

The ARDL($p$, $q$) model combines $p$ lags of $y$ with the
contemporaneous value and $q$ lags of each regressor $x$:
\[
  y_t = c
        + \sum_{i=1}^p \phi_i\, y_{t-i}
        + \sum_{j=0}^q \beta_j\, x_{t-j}
        + \varepsilon_t
\]

The \textbf{long-run multiplier} of $x$ on $y$ — the total accumulated
effect after all dynamic adjustments — is:
\[
  \theta = \frac{\sum_{j=0}^q \beta_j}{1 - \sum_{i=1}^p \phi_i}
\]
For the model to be stationary, $\left|\sum_{i=1}^p \phi_i\right| < 1$ is required.

\subsection*{Relationship with the VECM and the Pesaran test}

When $y$ and $x$ are integrated of order 1 ($I(1)$), the ARDL can be
rewritten in error correction form. Pesaran, Shin and Smith (2001)
propose the \emph{bounds test}: tests $H_0$ of no long-run relationship
via the $F$ statistic applied to the unrestricted ARDL, comparing with
tabulated critical values. The advantage over Johansen (ch.~\ref{cap:coint}) is that
the bounds test works regardless of whether $x$ is $I(0)$ or $I(1)$.

% ─────────────────────────────────────────────────────────────────────────────
\section{DGP Verification}

DGP: $y_t = 1 + 2x_t + \varepsilon_t$, no autocorrelation in $y$.
$N = 300$, $x \sim U(-\sqrt{3}, \sqrt{3})$, $\varepsilon \sim U(-0{,}5, 0{,}5)$.

\begin{lstlisting}[caption={AutoReg and ARDL}]
seed(42)
// ... (300 obs, generation of x and e)
generate df y = 1.0 + 2.0 * x + 0.5 * (rnormal() - 0.5)

// Pure AR(2): no true autocorrelation in y → coefs ~0
let m_ar   = autoreg(df, y, lags=2)
display m_ar

// AR(1) with deterministic trend
let m_ar_ct = autoreg(df, y, lags=1, trend="ct")
display m_ar_ct

// ARDL(1,1)
let m_ardl  = ardl(y ~ x, df, lags=1, xlags=1)
display m_ardl

// ARDL(2,2): check whether extra lags are significant
let m_ardl2 = ardl(y ~ x, df, lags=2, xlags=2)
display m_ardl2
\end{lstlisting}

\subsection*{Pure AR(2)}

\begin{lstlisting}[language={}, basicstyle=\ttfamily\small, frame=none,
  numbers=none, backgroundcolor=\color{codbg}]
================================= AutoReg(2) =================================
Observations:               298
R-squared:               0.0004
AIC:                   905.2221

-------------------------------- Coefficients --------------------------------
                        coef      std err            t        P>|t|
------------------------------------------------------------------------------
const                 0.9904       0.1022       9.6914       0.0000
y.L1                  0.0007       0.0580       0.0118       0.9906
y.L2                 -0.0209       0.0581      -0.3597       0.7193
==============================================================================
\end{lstlisting}

$R^2 = 0{,}0004$ — lags of $y$ explain nothing, as per the DGP
(no true autocorrelation). Coefficients \texttt{y.L1} and \texttt{y.L2}
do not differ from zero.

\subsection*{AR(1) with trend}

\begin{lstlisting}[language={}, basicstyle=\ttfamily\small, frame=none,
  numbers=none, backgroundcolor=\color{codbg}]
================================= AutoReg(1) =================================
Trend:                       ct

                        coef      std err            t        P>|t|
------------------------------------------------------------------------------
const                 0.9275       0.1394       6.6551       0.0000
trend                 0.0003       0.0007       0.4172       0.6768
y.L1                  0.0018       0.0581       0.0310       0.9753
==============================================================================
\end{lstlisting}

The deterministic trend ($p = 0{,}68$) and the lag ($p = 0{,}98$) are not
significant — consistent with a DGP without trend or persistence.

\subsection*{ARDL(1,1)}

\begin{lstlisting}[language={}, basicstyle=\ttfamily\small, frame=none,
  numbers=none, backgroundcolor=\color{codbg}]
================================= ARDL(1, 1) =================================
Observations:               299
R-squared:               0.9836
AIC:                  -318.6638

                        coef      std err            t        P>|t|
------------------------------------------------------------------------------
const                 1.0018       0.0594      16.8580       0.0000
y.L1                  0.0062       0.0583       0.1061       0.9155
x1                    1.9648       0.0148     133.1872       0.0000
x1.L1                -0.0078       0.1156      -0.0675       0.9462
==============================================================================
\end{lstlisting}

The model recovers the DGP accurately: \texttt{const} $\approx 1{,}0$,
\texttt{x1} $\approx 2{,}0$. The coefficient of \texttt{x1.L1} is not
significant ($p = 0{,}95$), confirming no lagged effect of $x$
in the DGP. The long-run multiplier is:
\[
  \theta = \frac{1{,}9648 + (-0{,}0078)}{1 - 0{,}0062} \approx 1{,}97
\]
close to the true value of 2{,}0.

\subsection*{ARDL(2,2)}

\begin{lstlisting}[language={}, basicstyle=\ttfamily\small, frame=none,
  numbers=none, backgroundcolor=\color{codbg}]
================================= ARDL(2, 2) =================================
Observations:               298
R-squared:               0.9838
AIC:                  -317.3857

                        coef      std err            t        P>|t|
------------------------------------------------------------------------------
const                 1.0727       0.0833      12.8843       0.0000
y.L1                  0.0013       0.0582       0.0220       0.9824
y.L2                 -0.0662       0.0582      -1.1385       0.2558
x1                    1.9628       0.0147     133.1320       0.0000
x1.L1                 0.0018       0.1154       0.0154       0.9877
x1.L2                 0.1264       0.1153       1.0968       0.2736
==============================================================================
\end{lstlisting}

The AIC of ARDL(2,2) ($-317{,}4$) is slightly \emph{higher} (worse) than that
of ARDL(1,1) ($-318{,}7$) — the extra lags do not compensate for the loss of
degrees of freedom. AIC/BIC selection favors the more parsimonious model.

% ─────────────────────────────────────────────────────────────────────────────
\section{Syntax}

\begin{lstlisting}
// AR(p) via OLS
autoreg(df, var, lags=1)                       // AR(1) with constant
autoreg(df, var, lags=p, trend="ct")           // AR(p) with trend

// ARDL(p,q)
ardl(y ~ x, df, lags=p, xlags=q)              // one regressor
ardl(gdp ~ inv + exp, df, lags=2, xlags=2)    // multiple regressors
\end{lstlisting}

\begin{center}
\begin{tabular}{lll}
\toprule
\textbf{Model} & \textbf{Command} & \textbf{Use case} \\
\midrule
AR($p$)      & \texttt{autoreg(df, y, lags=p)}       & Pure dynamics of $y$ \\
ARIMA($p,d,q$) & \texttt{arima(df, y, p=, d=, q=)}  & $I(1)$ with MA \\
ARDL($p,q$)  & \texttt{ardl(y{\raise.17ex\hbox{$\scriptstyle\sim$}}x, df, lags=p, xlags=q)} & $y$ and $x$ with lags \\
VECM         & \texttt{vecm(...)}                     & Multivariate cointegration \\
\bottomrule
\end{tabular}
\end{center}

\begin{nota}
  For $I(1)$ series, the ARDL can be rewritten in error correction form.
  The Pesaran et al.\ (2001) bounds test tests the existence of a
  long-run relationship without requiring all regressors to be $I(0)$
  or all $I(1)$ — an advantage over the Johansen test. The $F$ statistic
  of the test is accessible via \texttt{m\_ardl.bounds\_test(y, x)} in the
  Rust API of Greeners, but is not yet exposed as a DSL command.
\end{nota}
